% part: intuitionistic-logic
% chapter: soundness-completeness
% section: truth-lemma

\documentclass[../../../include/open-logic-section]{subfiles}

\begin{document}

\olfileid{int}{sc}{tru}

\olsection{સત્યતા સહાયક પ્રમેય}

નીચેનું સહાયક પ્રમેય કૅનોનિકલ નિદર્શમાં સત્ય હોવાને
કયા !!{formula}s અભાજ્ય ગણ~$\Delta$ના !!{element}s
છે તેની સાથે જોડે છે.

\begin{lem}\ollabel{lem:truth}
  જો $\Delta$ અભાજ્ય હોય, તો $\mSat{M(\Delta)}{!A}[\sigma]$
  ત્યારે અને માત્ર ત્યારે જ જ્યારે
  $\Delta(\sigma) \Proves !A$.
\end{lem}

\begin{proof}
  $!A$ પર આગમન કરીએ.
  \begin{enumerate}
    \item \indcase{!A}{\lfalse}{$\Delta(\sigma)$ અભાજ્ય છે,
      તેથી સુસંગત છે અને $\Delta(\sigma) \Proves/ \indfrm$.
      વ્યાખ્યા મુજબ $\mSat/{M(\Delta)}{\indfrm}[\sigma]$.}
      \item \indcase{!A}{p}{$\mSat{{}}{}$ની વ્યાખ્યા મુજબ
        $\mSat{M(\Delta)}{\indfrm}[\sigma]$ ત્યારે અને માત્ર
        ત્યારે જ જ્યારે $\sigma \in V(p)$, એટલે
        $\Delta(\sigma) \Proves \indfrm$.}
      \item \indcase!{!A}{\lnot !B}{}
      \item \indcase{!A}{!B \land
        !C}{$\mSat{M(\Delta)}{\indfrm}[\sigma]$ ત્યારે અને
        માત્ર ત્યારે જ જ્યારે
        $\mSat{M(\Delta)}{!B}[\sigma]$ અને
        $\mSat{M(\Delta)}{!C}[\sigma]$. આગમનની ધારણાથી
        $\mSat{M(\Delta)}{!B}[\sigma]$ ત્યારે અને માત્ર
        ત્યારે જ જ્યારે $\Delta(\sigma) \Proves
        !B$; $!C$ માટે પણ એ જ રીતે. હવે
        $\Delta(\sigma) \Proves !B$ અને
        $\Delta(\sigma) \Proves !C$ ત્યારે અને માત્ર ત્યારે
        જ જ્યારે $\Delta(\sigma)
        \Proves \indfrm$.}
      \item \indcase{!A}{!B \lor
        !C}{$\mSat{M(\Delta)}{\indfrm}[\sigma]$ ત્યારે અને
        માત્ર ત્યારે જ જ્યારે
        $\mSat{M(\Delta)}{!B}[\sigma]$ અથવા
        $\mSat{M(\Delta)}{!C}[\sigma]$. આગમનની ધારણાથી
        આ ત્યારે અને માત્ર ત્યારે જ જ્યારે
        $\Delta(\sigma) \Proves !B$ અથવા
        $\Delta(\sigma) \Proves !C$. હવે આ શરત
        $\Delta(\sigma) \Proves \indfrm$ સાથે સમકક્ષ
        છે તે બતાવવું છે. ડાબેથી જમણી દિશા સ્પષ્ટ છે.
        જમણેથી ડાબી દિશા $\Delta(\sigma)$ અભાજ્ય
        હોવાથી મળે છે.}
      \item \indcase{!A}{!B \lif !C}{પહેલાં ડાબેથી જમણી
        દિશાનું પ્રતિપક્ષી વિધાન. ધારો કે
        $\Delta(\sigma) \Proves/ !B
        \lif !C$. તો $\Delta(\sigma) \cup \{!B\} \Proves/
        !C$ પણ છે. કોઈક~$n$ માટે $\tuple{!B, !C}$ એ
        $\tuple{!B_n, !C_n}$ હોવાથી
        $\Delta(\sigma.n) = (\Delta(\sigma) \cup
        \{!B\})^*$, અને $\Delta(\sigma.n) \Proves !B$
        પરંતુ $\Delta(\sigma.n) \Proves/ !C$. આગમનની
        ધારણાથી $\mSat{M(\Delta)}{!B}[\sigma.n]$
        અને $\mSat/{M(\Delta)}{!C}[\sigma.n]$.
        $R\sigma(\sigma.n)$ હોવાથી આનો અર્થ
        $\mSat/{M(\Delta)}{\indfrm}[\sigma]$ થાય છે.

        હવે ધારો કે $\Delta(\sigma) \Proves !B \lif !C$
        અને $R\sigma\sigma'$. $\Delta(\sigma) \subseteq
        \Delta(\sigma')$ હોવાથી, જો
        $\Delta(\sigma') \Proves !B$ હોય, તો
        $\Delta(\sigma') \Proves !C$. બીજા શબ્દોમાં,
        $R\sigma\sigma'$ ધરાવતા દરેક $\sigma'$ માટે
        $\Delta(\sigma') \Proves/ !B$ અથવા
        $\Delta(\sigma') \Proves !C$. આગમનની ધારણાથી
        $R\sigma\sigma'$ હોય ત્યારે
        $\mSat/{M(\Delta)}{!B}[\sigma']$ અથવા
        $\mSat{M(\Delta)}{!C}[\sigma']$;
        એટલે $\mSat{M(\Delta)}{\indfrm}[\sigma]$.}
  \end{enumerate}
\end{proof}

\end{document}
